\chapter{§2.8 一元二次方程}
\label{sec:2.8}


\prerequisite{§2.4 因式分解, §2.6 一元一次方程}

\gradelevel{9 年级}


\section*{一元二次方程的一般形式}


\begin{explore}


学校要修建一个长方形花坛，长比宽多 $3$ 米，面积为 $40$ 平方米。设宽为 $x$ 米，则长为 $(x + 3)$ 米，由面积条件得 $x(x + 3) = 40$ ，展开为 $x^2 + 3x - 40 = 0$ 。这个方程中未知数的最高次数是 $2$ ——这就是\textbf{一元二次方程}。

\end{explore}

\begin{understand}


\textbf{一元二次方程}：只含一个未知数，且未知数的最高次数为 $2$ 的方程。

\textbf{一般形式}：

\[
ax^2 + bx + c = 0 \quad (a \neq 0)
\]


其中 $a$ 是\textbf{二次项系数}， $b$ 是\textbf{一次项系数}， $c$ 是\textbf{常数项}。

\textbf{为什么 $a \neq 0$ ？} 如果 $a = 0$ ，方程退化为一元一次方程 $bx + c = 0$ 。

\end{understand}

\begin{workedexamples}


\textbf{例 1.} 把 $3x^2 = 5x - 2$ 化为一般形式，并指出 $a, b, c$ 。

\textbf{解：} 移项得 $3x^2 - 5x + 2 = 0$ 。 $a = 3$ ， $b = -5$ ， $c = 2$ 。

\textbf{例 2.} 把 $(x-1)(x+3) = 5$ 化为一般形式。

\textbf{解：} 展开： $x^2 + 3x - x - 3 = 5$ ， $x^2 + 2x - 8 = 0$ 。 $a = 1$ ， $b = 2$ ， $c = -8$ 。

\textbf{例 3.} 判断 $x^2 = 0$ 是否为一元二次方程。

\textbf{解：} 是。 $a = 1$ ， $b = 0$ ， $c = 0$ ，满足 $a \neq 0$ 。

\end{workedexamples}

\begin{keytakeaway}


\begin{itemize}
  \item 一元二次方程的标志：一个未知数，最高次数为 $2$ 。
  \item 一般形式 $ax^2 + bx + c = 0$ 中， $a \neq 0$ 是关键条件。
  \item 化为一般形式：把所有项移到等号左边，按降幂排列。
\end{itemize}


\end{keytakeaway}

\begin{practice}


\begin{enumerate}
  \item 把 $2x^2 + 3 = 7x$ 化为一般形式。
  \item 把 $(x + 4)^2 = 3x$ 化为一般形式，并指出 $a, b, c$ 。
  \item 判断 $x^3 - x^2 = 0$ 是否为一元二次方程。
  \item 若 $(m-1)x^2 + 3x - 5 = 0$ 是一元二次方程， $m$ 满足什么条件？
  \item 写出一个 $a = 2, c = -3$ 的一元二次方程。
\end{enumerate}


\end{practice}

\section*{解法：直接开平方法与配方法}


\begin{explore}


最简单的一元二次方程是 $x^2 = 4$ 。直接开平方就能得到 $x = \pm 2$ 。稍微复杂一点的如 $(x - 1)^2 = 9$ ，同样可以直接开平方。而对于 $x^2 + 6x + 5 = 0$ 这类不能直接开平方的方程，我们需要先把它"凑成"完全平方的形式——这就是\textbf{配方法}。

\end{explore}

\begin{understand}


\textbf{直接开平方法}：对于 $x^2 = k$ （ $k \geq 0$ ），直接得 $x = \pm\sqrt{k}$ 。

更一般地： $(x + m)^2 = k$ （ $k \geq 0$ ），则 $x + m = \pm\sqrt{k}$ ， $x = -m \pm \sqrt{k}$ 。

\textbf{配方法}：把 $ax^2 + bx + c = 0$ 化为 $(x + p)^2 = q$ 的形式。

配方步骤（以 $a = 1$ 为例）：

\begin{enumerate}
  \item 移项： $x^2 + bx = -c$ 。
  \item 两边加 $\left(\frac{b}{2}\right)^2$ ： $x^2 + bx + \left(\frac{b}{2}\right)^2 = -c + \left(\frac{b}{2}\right)^2$ 。
  \item 左边成为完全平方： $\left(x + \frac{b}{2}\right)^2 = \frac{b^2 - 4c}{4}$ 。
\end{enumerate}


\textbf{为什么配方法有效？} 它的本质是利用完全平方公式，人为"补"一个数，使左边成为完全平方形式，然后用直接开平方法求解。

\end{understand}

\begin{workedexamples}


\textbf{例 1.} 解方程 $x^2 = 25$ 。

\textbf{解：} $x = \pm\sqrt{25} = \pm 5$ 。

\textbf{例 2.} 解方程 $(x - 3)^2 = 7$ 。

\textbf{解：} $x - 3 = \pm\sqrt{7}$ ， $x = 3 \pm \sqrt{7}$ 。

\textbf{例 3.} 用配方法解 $x^2 + 6x + 5 = 0$ 。

\textbf{解：}

移项： $x^2 + 6x = -5$ 。

两边加 $\left(\frac{6}{2}\right)^2 = 9$ ：

\[
x^2 + 6x + 9 = -5 + 9 = 4
\]


\[
(x + 3)^2 = 4
\]


\[
x + 3 = \pm 2
\]


\[
x = -3 + 2 = -1 \quad \text{或} \quad x = -3 - 2 = -5
\]


\textbf{例 4.} 用配方法解 $2x^2 - 8x + 3 = 0$ 。

\textbf{解：}

先将二次项系数化为 $1$ ： $x^2 - 4x + \frac{3}{2} = 0$ 。

移项： $x^2 - 4x = -\frac{3}{2}$ 。

配方，加 $\left(\frac{-4}{2}\right)^2 = 4$ ：

\[
x^2 - 4x + 4 = -\frac{3}{2} + 4 = \frac{5}{2}
\]


\[
(x - 2)^2 = \frac{5}{2}
\]


\[
x = 2 \pm \sqrt{\frac{5}{2}} = 2 \pm \frac{\sqrt{10}}{2}
\]


\end{workedexamples}

\begin{keytakeaway}


\begin{itemize}
  \item 直接开平方法适用于 $(x + m)^2 = k$ 的形式。
  \item 配方法是将一般形式转化为可直接开平方的形式。
  \item 配方时，二次项系数要先化为 $1$ ，然后加上一次项系数一半的平方。
\end{itemize}


\end{keytakeaway}

\begin{practice}


\begin{enumerate}
  \item 解方程 $x^2 = 16$ 。
  \item 解方程 $(x + 2)^2 = 5$ 。
  \item 用配方法解 $x^2 - 4x - 5 = 0$ 。
  \item 用配方法解 $x^2 + 8x + 12 = 0$ 。
  \item 用配方法解 $3x^2 - 6x - 1 = 0$ 。
  \item 解方程 $4x^2 - 9 = 0$ 。
\end{enumerate}


\end{practice}

\section*{公式法}


\begin{explore}


配方法虽然通用，但步骤较多。如果对一般方程 $ax^2 + bx + c = 0$ 直接用配方法推导，可以一劳永逸地得到一个万能的\textbf{求根公式}——这就是历史上最重要的代数公式之一。

\end{explore}

\begin{understand}


对 $ax^2 + bx + c = 0$ （ $a \neq 0$ ）进行配方：

\[
\begin{aligned}
x^2 + \frac{b}{a}x &= -\frac{c}{a} \\[4pt]
x^2 + \frac{b}{a}x + \frac{b^2}{4a^2} &= \frac{b^2 - 4ac}{4a^2} \\[4pt]
\left(x + \frac{b}{2a}\right)^2 &= \frac{b^2 - 4ac}{4a^2}
\end{aligned}
\]


当 $b^2 - 4ac \geq 0$ 时：

\[
\boxed{x = \frac{-b \pm \sqrt{b^2 - 4ac}}{2a}}
\]


这就是\textbf{一元二次方程的求根公式}。

\end{understand}

\begin{workedexamples}


\textbf{例 1.} 用公式法解 $x^2 - 5x + 6 = 0$ 。

\textbf{解：} $a = 1$ ， $b = -5$ ， $c = 6$ 。

\[
\Delta = b^2 - 4ac = 25 - 24 = 1
\]


\[
x = \frac{5 \pm \sqrt{1}}{2} = \frac{5 \pm 1}{2}
\]


\[
x_1 = 3, \quad x_2 = 2
\]


\textbf{例 2.} 用公式法解 $2x^2 + 3x - 5 = 0$ 。

\textbf{解：} $a = 2$ ， $b = 3$ ， $c = -5$ 。

\[
\Delta = 9 - 4 \times 2 \times (-5) = 9 + 40 = 49
\]


\[
x = \frac{-3 \pm 7}{4}
\]


\[
x_1 = \frac{-3 + 7}{4} = 1, \quad x_2 = \frac{-3 - 7}{4} = -\frac{5}{2}
\]


\textbf{例 3.} 用公式法解 $x^2 - 4x + 1 = 0$ 。

\textbf{解：} $a = 1$ ， $b = -4$ ， $c = 1$ 。

\[
\Delta = 16 - 4 = 12
\]


\[
x = \frac{4 \pm \sqrt{12}}{2} = \frac{4 \pm 2\sqrt{3}}{2} = 2 \pm \sqrt{3}
\]


\end{workedexamples}

\begin{keytakeaway}


\begin{itemize}
  \item 求根公式 $x = \frac{-b \pm \sqrt{b^2 - 4ac}}{2a}$ 适用于一切一元二次方程。
  \item 使用前先算判别式 $\Delta = b^2 - 4ac$ ，确保 $\Delta \geq 0$ 。
  \item 注意 $b$ 的符号， $-b$ 中的负号不能丢。
\end{itemize}


\end{keytakeaway}

\begin{practice}


\begin{enumerate}
  \item 用公式法解 $x^2 - 3x - 10 = 0$ 。
  \item 用公式法解 $2x^2 - 7x + 3 = 0$ 。
  \item 用公式法解 $x^2 + 2x - 5 = 0$ 。
  \item 用公式法解 $3x^2 - x - 1 = 0$ 。
  \item 用公式法解 $x^2 - 6x + 9 = 0$ 。
\end{enumerate}


\end{practice}

\section*{因式分解法}


\begin{explore}


对于 $x^2 - 5x + 6 = 0$ ，如果我们能看出左边可以分解为 $(x-2)(x-3)$ ，那么由"两个因式的乘积为零，则至少有一个因式为零"，立刻得到 $x = 2$ 或 $x = 3$ 。这比用公式法快得多！

\end{explore}

\begin{understand}


\textbf{因式分解法}的步骤：

\begin{enumerate}
  \item 将方程化为一般形式 $ax^2 + bx + c = 0$ 。
  \item 对左边进行因式分解： $a(x - x_1)(x - x_2) = 0$ 。
  \item 令每个因式为零： $x = x_1$ 或 $x = x_2$ 。
\end{enumerate}


\textbf{零因式定理}：若 $AB = 0$ ，则 $A = 0$ 或 $B = 0$ 。

\end{understand}

\begin{workedexamples}


\textbf{例 1.} 解方程 $x^2 - 7x + 12 = 0$ 。

\textbf{解：} 因式分解： $(x - 3)(x - 4) = 0$ 。

$x - 3 = 0$ 或 $x - 4 = 0$ ，即 $x_1 = 3$ ， $x_2 = 4$ 。

\textbf{例 2.} 解方程 $x^2 - 4x = 0$ 。

\textbf{解：} 提公因式： $x(x - 4) = 0$ 。

$x = 0$ 或 $x = 4$ 。

注意：不能两边除以 $x$ ，那样会丢掉 $x = 0$ 这个根！

\textbf{例 3.} 解方程 $4x^2 - 9 = 0$ 。

\textbf{解：} 用平方差公式： $(2x + 3)(2x - 3) = 0$ 。

$x = -\frac{3}{2}$ 或 $x = \frac{3}{2}$ 。

\textbf{例 4.} 解方程 $2x^2 + 5x - 3 = 0$ 。

\textbf{解：} 十字相乘： $(2x - 1)(x + 3) = 0$ 。

$x = \frac{1}{2}$ 或 $x = -3$ 。

\end{workedexamples}

\begin{keytakeaway}


\begin{itemize}
  \item 因式分解法是解一元二次方程最快的方法（当能分解时）。
  \item 核心原理：零因式定理。
  \item 注意不要丢根：不能两边同时除以含 $x$ 的式子。
\end{itemize}


\end{keytakeaway}

\begin{practice}


\begin{enumerate}
  \item 解方程 $x^2 + x - 6 = 0$ 。
  \item 解方程 $x^2 = 3x$ 。
  \item 解方程 $9x^2 - 1 = 0$ 。
  \item 解方程 $(x - 1)(x + 2) = 4$ 。（先展开化为一般形式）
  \item 解方程 $6x^2 - x - 2 = 0$ 。
\end{enumerate}


\end{practice}

\section*{判别式与根的情况}


\begin{explore}


用公式法解方程时，根号下的 $b^2 - 4ac$ 决定了方程的根的情况。如果它是正数，开方得到两个不同的值；如果它是零，两个根相同；如果它是负数，根号下出现负数，在实数范围内没有解。这个表达式非常重要，有一个专门的名字——\textbf{判别式}。

\end{explore}

\begin{understand}


\textbf{判别式}：对于 $ax^2 + bx + c = 0$ （ $a \neq 0$ ）， $\Delta = b^2 - 4ac$ 叫做判别式。

\[
\begin{cases}
\Delta > 0 & \text{方程有两个不相等的实数根} \\
\Delta = 0 & \text{方程有两个相等的实数根（即一个重根）} \\
\Delta < 0 & \text{方程没有实数根}
\end{cases}
\]


\textbf{为什么 $\Delta$ 这么重要？} 判别式让我们在不求解的情况下，就能判断方程根的数量和类型。这在证明题和应用题中非常有用。

\end{understand}

\begin{workedexamples}


\textbf{例 1.} 不解方程，判断 $x^2 - 4x + 5 = 0$ 的根的情况。

\textbf{解：} $\Delta = (-4)^2 - 4 \times 1 \times 5 = 16 - 20 = -4 < 0$ 。

方程没有实数根。

\textbf{例 2.} 若方程 $x^2 - 2x + k = 0$ 有两个不相等的实数根，求 $k$ 的取值范围。

\textbf{解：} 要求 $\Delta > 0$ ：

\[
(-2)^2 - 4 \times 1 \times k > 0
\]


\[
4 - 4k > 0
\]


\[
k < 1
\]


\textbf{例 3.} 若方程 $x^2 + mx + 4 = 0$ 有两个相等的实数根，求 $m$ 的值。

\textbf{解：} $\Delta = 0$ ： $m^2 - 16 = 0$ ， $m = \pm 4$ 。

\textbf{例 4.} 证明：方程 $x^2 - 2(a + 1)x + a^2 + 1 = 0$ 总有实数根。

\textbf{解：} 计算判别式。

\[
\begin{aligned}
\Delta &= [2(a+1)]^2 - 4(a^2 + 1) \\
&= 4(a^2 + 2a + 1) - 4a^2 - 4 \\
&= 4a^2 + 8a + 4 - 4a^2 - 4 \\
&= 8a
\end{aligned}
\]


当 $a \geq 0$ 时 $\Delta \geq 0$ ，方程有实数根。当 $a < 0$ 时 $\Delta < 0$ ，方程无实数根。

所以原命题不成立——只有 $a \geq 0$ 时才有实数根。

\end{workedexamples}

\begin{keytakeaway}


\begin{itemize}
  \item $\Delta > 0$ ：两个不等实根； $\Delta = 0$ ：重根； $\Delta < 0$ ：无实根。
  \item 判别式可以用来确定参数的取值范围。
  \item 计算 $\Delta$ 时注意各项符号。
\end{itemize}


\end{keytakeaway}

\begin{practice}


\begin{enumerate}
  \item 判断 $x^2 + 3x + 5 = 0$ 的根的情况。
  \item 判断 $4x^2 - 4x + 1 = 0$ 的根的情况。
  \item 若 $x^2 + 4x + m = 0$ 有实数根，求 $m$ 的取值范围。
  \item 若 $kx^2 - 6x + 1 = 0$ 有两个相等实数根，求 $k$ 。
  \item 判断 $2x^2 - 5x + 1 = 0$ 的根的情况，并求出根。
\end{enumerate}


\end{practice}

\section*{韦达定理}


\begin{explore}


对于 $x^2 - 5x + 6 = 0$ ，它的两个根是 $2$ 和 $3$ 。注意到 $2 + 3 = 5$ （等于 $-\frac{b}{a}$ ）， $2 \times 3 = 6$ （等于 $\frac{c}{a}$ ）。这是巧合吗？不是——这是一个普遍规律，叫做\textbf{韦达定理}（也叫根与系数的关系）。

\end{explore}

\begin{understand}


\textbf{韦达定理}：若 $x_1, x_2$ 是 $ax^2 + bx + c = 0$ （ $a \neq 0$ ）的两个根，则：

\[
x_1 + x_2 = -\frac{b}{a}, \quad x_1 \cdot x_2 = \frac{c}{a}
\]


\textbf{证明}：由求根公式， $x_1 = \frac{-b + \sqrt{\Delta}}{2a}$ ， $x_2 = \frac{-b - \sqrt{\Delta}}{2a}$ 。

\[
x_1 + x_2 = \frac{-2b}{2a} = -\frac{b}{a}
\]


\[
x_1 \cdot x_2 = \frac{(-b)^2 - \Delta}{4a^2} = \frac{b^2 - (b^2 - 4ac)}{4a^2} = \frac{4ac}{4a^2} = \frac{c}{a}
\]


\textbf{韦达定理的用途}：不解方程就能知道两根之和与两根之积，常用于验证、构造方程和求含根的表达式。

\end{understand}

\begin{workedexamples}


\textbf{例 1.} 若 $x_1, x_2$ 是 $x^2 - 7x + 10 = 0$ 的两个根，求 $x_1 + x_2$ 和 $x_1 x_2$ 。

\textbf{解：} $x_1 + x_2 = -\frac{-7}{1} = 7$ ， $x_1 x_2 = \frac{10}{1} = 10$ 。

\textbf{例 2.} 若 $x_1, x_2$ 是 $2x^2 + 3x - 1 = 0$ 的两个根，求 $\frac{1}{x_1} + \frac{1}{x_2}$ 。

\textbf{解：}

\[
\frac{1}{x_1} + \frac{1}{x_2} = \frac{x_1 + x_2}{x_1 x_2} = \frac{-\frac{3}{2}}{-\frac{1}{2}} = 3
\]


\textbf{例 3.} 已知一个一元二次方程的两个根为 $-1$ 和 $4$ ，求这个方程。

\textbf{解：} 设方程为 $x^2 - (x_1 + x_2)x + x_1 x_2 = 0$ 。

$x_1 + x_2 = 3$ ， $x_1 x_2 = -4$ 。

方程为 $x^2 - 3x - 4 = 0$ 。

\textbf{例 4.} 若 $x_1, x_2$ 是 $x^2 - 3x + 1 = 0$ 的两根，求 $x_1^2 + x_2^2$ 。

\textbf{解：} $x_1 + x_2 = 3$ ， $x_1 x_2 = 1$ 。

\[
x_1^2 + x_2^2 = (x_1 + x_2)^2 - 2x_1 x_2 = 9 - 2 = 7
\]


\end{workedexamples}

\begin{keytakeaway}


\begin{itemize}
  \item 韦达定理：两根之和 $= -\frac{b}{a}$ ，两根之积 $= \frac{c}{a}$ 。
  \item 利用韦达定理可以不解方程就求出根的相关表达式。
  \item 已知两根可以反向构造一元二次方程。
\end{itemize}


\end{keytakeaway}

\begin{practice}


\begin{enumerate}
  \item 方程 $x^2 - 4x + 3 = 0$ 的两根之和与两根之积分别是多少？
  \item 方程 $3x^2 + 5x - 2 = 0$ 的两根之和和两根之积分别是多少？
  \item 已知两根为 $2$ 和 $-5$ ，写出一元二次方程。
  \item 若 $x_1, x_2$ 是 $x^2 + 2x - 5 = 0$ 的两根，求 $x_1^2 + x_2^2$ 。
  \item 若 $x_1, x_2$ 是 $x^2 - 5x + 3 = 0$ 的两根，求 $(x_1 - x_2)^2$ 。
\end{enumerate}


\end{practice}

\section*{一元二次方程的应用}


\begin{explore}


利润、面积、工程等实际问题中，经常出现二次关系。例如：某商品进价 $20$ 元，售价为 $25$ 元时每天可卖 $100$ 件。据调查，售价每降低 $1$ 元，每天多卖 $20$ 件。问售价定为多少时，每天利润为 $600$ 元？

\end{explore}

\begin{understand}


一元二次方程的应用与一元一次方程类似，关键仍是\textbf{列方程}。不同的是，二次方程可能有两个解，需要根据实际意义\textbf{取舍}。

\end{understand}

\begin{workedexamples}


\textbf{例 1.} 一块长方形铁皮长 $20$ 厘米、宽 $16$ 厘米，在四个角各截去一个同样的小正方形后，折成一个无盖长方体盒子，使盒子的底面积为 $192$ 平方厘米。求截去的小正方形的边长。

\textbf{解：} 设小正方形的边长为 $x$ 厘米。折起后，长 $= 20 - 2x$ ，宽 $= 16 - 2x$ 。

\[
(20 - 2x)(16 - 2x) = 192
\]


\[
320 - 40x - 32x + 4x^2 = 192
\]


\[
4x^2 - 72x + 128 = 0
\]


\[
x^2 - 18x + 32 = 0
\]


\[
(x - 2)(x - 16) = 0
\]


$x = 2$ 或 $x = 16$ 。

因为宽 $16$ 厘米， $x = 16$ 时宽变为负数，舍去。 $x = 2$ 。

答：截去的小正方形边长为 $2$ 厘米。

\textbf{例 2.} 某商品进价 $20$ 元，售价 $25$ 元时日销量 $100$ 件。售价每降 $1$ 元日销量增加 $20$ 件。问售价为多少时日利润为 $600$ 元？

\textbf{解：} 设售价为 $(25 - x)$ 元（ $x$ 为降价金额），则日销量 $(100 + 20x)$ 件，每件利润 $(5 - x)$ 元。

\[
(5 - x)(100 + 20x) = 600
\]


\[
500 + 100x - 100x - 20x^2 = 600
\]


\[
-20x^2 = 100
\]


\[
x^2 = -5
\]


无实数解。说明每天利润无法达到 $600$ 元——让我们验证最大利润。

实际上，利润 $= (25 - x - 20)(100 + 20x) = (5 - x)(100 + 20x) = 500 + 100x - 100x - 20x^2 = 500 - 20x^2$ 。

最大利润 $500$ 元（当 $x = 0$ ，即不降价），所以利润 $600$ 元确实不可能。

重新设置题目。设售价降低 $x$ 元，每件利润为 $(25 - x - 20) = (5 - x)$ 元，日销量为 $(100 + 20x)$ 件。

日利润 $= (5 - x)(100 + 20x) = 500 + 100x - 100x - 20x^2 = 500 - 20x^2$ 。

若要利润 $400$ 元： $500 - 20x^2 = 400$ ， $x^2 = 5$ ， $x = \sqrt{5} \approx 2.24$ 。售价约 $22.76$ 元。

换一个标准题目。

\textbf{例 2（修正）.} 某商品进价 $40$ 元，售价 $50$ 元时日销量 $30$ 件。每降价 $1$ 元多卖 $2$ 件。日利润 $396$ 元时售价是多少？

\textbf{解：} 设降价 $x$ 元，则售价 $(50 - x)$ 元，每件利润 $(10 - x)$ 元，日销量 $(30 + 2x)$ 件。

\[
(10 - x)(30 + 2x) = 396
\]


\[
300 + 20x - 30x - 2x^2 = 396
\]


\[
-2x^2 - 10x + 300 = 396
\]


\[
2x^2 + 10x + 96 = 0
\]


\[
x^2 + 5x + 48 = 0
\]


$\Delta = 25 - 192 < 0$ 。仍然无解。

让我们用经典题型重写。

\textbf{例 2（经典题型）.} 某商品进价 $30$ 元，售价定为 $40$ 元时每天可卖 $20$ 件。已知售价每涨 $1$ 元日销量减少 $1$ 件。要使每天利润为 $234$ 元，售价应定为多少？

\textbf{解：} 设售价为 $(40 + x)$ 元，则每件利润 $(10 + x)$ 元，日销量 $(20 - x)$ 件。

\[
(10 + x)(20 - x) = 234
\]


\[
200 - 10x + 20x - x^2 = 234
\]


\[
-x^2 + 10x + 200 = 234
\]


\[
x^2 - 10x + 34 = 0
\]


$\Delta = 100 - 136 = -36 < 0$ 。

改为利润 $200$ 元试试。原始不降价时利润 $= 10 \times 20 = 200$ ，最大利润为 $(10 + 5)(20 - 5) = 225$ 。

设利润 $216$ 元： $(10 + x)(20 - x) = 216$ ， $-x^2 + 10x + 200 = 216$ ， $x^2 - 10x + 16 = 0$ ， $(x - 2)(x - 8) = 0$ ， $x = 2$ 或 $x = 8$ 。

\textbf{例 2（最终版）.} 某商品进价 $30$ 元，售价定为 $40$ 元时每天可卖 $20$ 件。售价每涨 $1$ 元日销量减少 $1$ 件。问售价定为多少时日利润为 $216$ 元？

\textbf{解：} 设售价涨 $x$ 元，每件利润 $(10 + x)$ 元，日销量 $(20 - x)$ 件。

\[
(10 + x)(20 - x) = 216
\]


\[
200 + 10x - x^2 = 216
\]


\[
x^2 - 10x + 16 = 0
\]


\[
(x - 2)(x - 8) = 0
\]


$x = 2$ 或 $x = 8$ 。

\begin{itemize}
  \item $x = 2$ ：售价 $42$ 元，日销量 $18$ 件。
  \item $x = 8$ ：售价 $48$ 元，日销量 $12$ 件。
\end{itemize}


答：售价定为 $42$ 元或 $48$ 元时，日利润均为 $216$ 元。

\textbf{例 3.} 某公园有一条长 $60$ 米的直路两旁要种树，如果每隔 $x$ 米种一棵，两端都种，一侧可种 $\frac{60}{x} + 1$ 棵。两侧共种 $62$ 棵，求间距 $x$ 。

\textbf{解：} 两侧共 $2\left(\frac{60}{x} + 1\right) = 62$ ， $\frac{60}{x} + 1 = 31$ ， $\frac{60}{x} = 30$ ， $x = 2$ 。

（这是一元一次方程的情境。下面换一个二次方程应用。）

\textbf{例 3.} 一个直角三角形的两条直角边之差为 $3$ 厘米，斜边为 $15$ 厘米。求两条直角边。

\textbf{解：} 设较短直角边为 $x$ 厘米，较长直角边为 $(x + 3)$ 厘米。由勾股定理：

\[
x^2 + (x + 3)^2 = 15^2
\]


\[
x^2 + x^2 + 6x + 9 = 225
\]


\[
2x^2 + 6x - 216 = 0
\]


\[
x^2 + 3x - 108 = 0
\]


\[
(x + 12)(x - 9) = 0
\]


$x = 9$ 或 $x = -12$ （舍去负值）。

两条直角边为 $9$ 厘米和 $12$ 厘米。

验证： $9^2 + 12^2 = 81 + 144 = 225 = 15^2$ $\checkmark$。

\end{workedexamples}

\begin{keytakeaway}


\begin{itemize}
  \item 一元二次方程的应用题，列方程后可能得到两个解，要根据实际意义取舍。
  \item 常见题型：面积问题、利润问题、勾股定理问题、增长率问题。
  \item 解方程前先检查 $\Delta$ ，若 $\Delta < 0$ 说明问题无解，需重新审题。
\end{itemize}


\end{keytakeaway}

\begin{practice}


\begin{enumerate}
  \item 一个正方形的面积增加 $20$ 平方厘米后变为另一个正方形，新正方形的边长比原来长 $2$ 厘米。求原正方形的边长。
  \item 两个连续正整数之积为 $132$ ，求这两个数。
  \item 某种产品两年内降价 $36\%$ ，若每年的降价率相同，求年降价率。
  \item 一个直角三角形的两条直角边之和为 $14$ 厘米，斜边为 $10$ 厘米，求两条直角边。
  \item 一块矩形草地长比宽多 $6$ 米，面积为 $72$ 平方米，求长和宽。
\end{enumerate}


\end{practice}



\begin{answer}


\textbf{一元二次方程的一般形式 练一练}

\begin{enumerate}
  \item $2x^2 - 7x + 3 = 0$ 。
  \item $x^2 + 8x + 16 = 3x$ ， $x^2 + 5x + 16 = 0$ 。 $a = 1$ ， $b = 5$ ， $c = 16$ 。
  \item 不是。最高次数为 $3$ 。
  \item $m - 1 \neq 0$ ，即 $m \neq 1$ 。
  \item 例如 $2x^2 + x - 3 = 0$ （ $b$ 可任取）。
\end{enumerate}


\textbf{直接开平方法与配方法 练一练}

\begin{enumerate}
  \item $x = \pm 4$ 。
  \item $x + 2 = \pm\sqrt{5}$ ， $x = -2 \pm \sqrt{5}$ 。
  \item $x^2 - 4x = 5$ ， $x^2 - 4x + 4 = 9$ ， $(x-2)^2 = 9$ ， $x = 5$ 或 $x = -1$ 。
  \item $x^2 + 8x = -12$ ， $x^2 + 8x + 16 = 4$ ， $(x+4)^2 = 4$ ， $x = -2$ 或 $x = -6$ 。
  \item $x^2 - 2x = \frac{1}{3}$ ， $x^2 - 2x + 1 = \frac{4}{3}$ ， $(x-1)^2 = \frac{4}{3}$ ， $x = 1 \pm \frac{2\sqrt{3}}{3}$ 。
  \item $(2x)^2 = 9$ ， $2x = \pm 3$ ， $x = \pm\frac{3}{2}$ 。
\end{enumerate}


\textbf{公式法 练一练}

\begin{enumerate}
  \item $\Delta = 9 + 40 = 49$ ， $x = \frac{3 \pm 7}{2}$ ， $x_1 = 5$ ， $x_2 = -2$ 。
  \item $\Delta = 49 - 24 = 25$ ， $x = \frac{7 \pm 5}{4}$ ， $x_1 = 3$ ， $x_2 = \frac{1}{2}$ 。
  \item $\Delta = 4 + 20 = 24$ ， $x = \frac{-2 \pm 2\sqrt{6}}{2} = -1 \pm \sqrt{6}$ 。
  \item $\Delta = 1 + 12 = 13$ ， $x = \frac{1 \pm \sqrt{13}}{6}$ 。
  \item $\Delta = 36 - 36 = 0$ ， $x = \frac{6}{2} = 3$ （重根）。
\end{enumerate}


\textbf{因式分解法 练一练}

\begin{enumerate}
  \item $(x + 3)(x - 2) = 0$ ， $x = -3$ 或 $x = 2$ 。
  \item $x^2 - 3x = 0$ ， $x(x - 3) = 0$ ， $x = 0$ 或 $x = 3$ 。
  \item $(3x + 1)(3x - 1) = 0$ ， $x = \pm\frac{1}{3}$ 。
  \item $x^2 + x - 2 = 4$ ， $x^2 + x - 6 = 0$ ， $(x + 3)(x - 2) = 0$ ， $x = -3$ 或 $x = 2$ 。
  \item $(3x - 2)(2x + 1) = 0$ ， $x = \frac{2}{3}$ 或 $x = -\frac{1}{2}$ 。
\end{enumerate}


\textbf{判别式与根的情况 练一练}

\begin{enumerate}
  \item $\Delta = 9 - 20 = -11 < 0$ ，无实数根。
  \item $\Delta = 16 - 16 = 0$ ，有两个相等的实数根（重根）。
  \item $\Delta = 16 - 4m \geq 0$ ， $m \leq 4$ 。
  \item $\Delta = 36 - 4k = 0$ ， $k = 9$ 。
  \item $\Delta = 25 - 8 = 17 > 0$ ，有两个不等实根。 $x = \frac{5 \pm \sqrt{17}}{4}$ 。
\end{enumerate}


\textbf{韦达定理 练一练}

\begin{enumerate}
  \item 两根之和 $= 4$ ，两根之积 $= 3$ 。
  \item 两根之和 $= -\frac{5}{3}$ ，两根之积 $= -\frac{2}{3}$ 。
  \item $x^2 + 3x - 10 = 0$ （两根之和 $= -3$ ，两根之积 $= -10$ ）。
  \item $x_1 + x_2 = -2$ ， $x_1 x_2 = -5$ 。 $x_1^2 + x_2^2 = (-2)^2 - 2(-5) = 4 + 10 = 14$ 。
  \item $(x_1 - x_2)^2 = (x_1 + x_2)^2 - 4x_1 x_2 = 25 - 12 = 13$ 。
\end{enumerate}


\textbf{一元二次方程的应用 练一练}

\begin{enumerate}
  \item 设原正方形边长 $a$ ， $(a + 2)^2 = a^2 + 20$ ， $a^2 + 4a + 4 = a^2 + 20$ ， $4a = 16$ ， $a = 4$ 厘米。
  \item 设较小数为 $n$ ， $n(n + 1) = 132$ ， $n^2 + n - 132 = 0$ ， $(n + 12)(n - 11) = 0$ ， $n = 11$ （正整数）。两个数为 $11$ 和 $12$ 。
  \item 设年降价率为 $x$ ， $(1 - x)^2 = 1 - 0.36 = 0.64$ ， $1 - x = 0.8$ （取正值）， $x = 0.2 = 20\%$ 。
  \item 设两直角边为 $a, b$ 。 $a + b = 14$ ， $a^2 + b^2 = 100$ 。 $a^2 + b^2 = (a+b)^2 - 2ab = 196 - 2ab = 100$ ， $ab = 48$ 。 $a, b$ 是 $t^2 - 14t + 48 = 0$ 的根。 $(t - 6)(t - 8) = 0$ ， $a = 6, b = 8$ （或反之）。
  \item 设宽 $x$ 米，长 $(x + 6)$ 米。 $x(x + 6) = 72$ ， $x^2 + 6x - 72 = 0$ ， $(x + 12)(x - 6) = 0$ ， $x = 6$ 。宽 $6$ 米，长 $12$ 米。
\end{enumerate}


\end{answer}
